%##%\FINALFIX{\newpage}
\subsection{nml\_initialization}
\label{SEC:nml_initialization}
\inp*[nmlInitialization]{}[nml\_initialization][]

This \hyperref[SEC:User_Input]{namelist block}
is used to specify solution initialization.

\begin{nmlblocklst}[
,label={LST:nml_initialization}
,caption={Namelist nml\_initialization.}
]
&nml_initialization
  Initialization                = 1
  Jet_Init_Func                 = 2
  Load_Restart_File             = F
  Restart_File                  = "out.rst"
  Initial_Solution_File         = ""
  Initial_Solution_Format       = -99        ! -99 = Unknown
/
\end{nmlblocklst}
% CFL_Reduction_For_Nan_Restart = 2.0

The \key*[key:InitializationSequence]{initialization sequence}[Initialization Sequence]
is defined with the following precedence:
\begin{itemize}
%
\item
If the file specified by \inp[nmlInitialization::InitialSolutionFile]{Initial\_Solution\_File}
exists and \inp[nmlInitialization::InitialSolutionFormat]{Initial\_Solution\_Format} is set to
a valid value, then the solution data read from the specified file will be
used as the initialization for a new simulation.
%
Note that \inptext{Initial\_Solution\_File} and \inptext{Initial\_Solution\_Format} can be
set to read an externally generated solution or a GFR restart file.
%
The iteration and solution time will be set to zero.
%
\item
If the file specified by \inp[nmlInitialization::RestartFile]{Restart\_File}
exists and
\inp[nmlInitialization::LoadRestartFile]{Load\_Restart\_File}=\texttt{T},
then the solution will be initialized using the data read from
the specified restart file.
The iteration and solution time will continue from that in the restart file.
%
\item
Otherwise, the initial solution will be defined based on
the \inp[nmlInitialization::Initialization]{Initialization} option.
The iteration and solution time will be set to zero.
\end{itemize}

\noindent\NOTE
To continue a simulation from a restart file while preserving the iteration
and solution time, one should set
  \inptext{Initial\_Solution\_File}=\texttt{""},
  \inptext{Load\_Restart\_File}=\texttt{T}, and
  \inptext{Restart\_File} to the name of the input restart file.
Many users forget to make these changes, and their simulations start over
from the previously specified initialization rather than continuing the
simulation.


\devNotes{
%
The files specified by \inptext{Initial\_Solution\_File} and \inptext{Restart\_File}
need to be checked for validity as part of the all the error checks after the
namelist is read.
%
Currently, I think it goes all the way to trying to read the specified restart
file before finding out if it does not exist, and at that stage it is not
immediately clear if it is \inptext{Initial\_Solution\_File} or \inptext{Restart\_File}
that is responsible.
%
Earlier error checking can allow for a more graceful abort with a clearer
message like
%
\begin{InputExample}
  \item[\quoterm{Tried to read \inptext{Initial\_Solution\_File}, but file does not
  exist!}]
\end{InputExample}
%
immediately after reading the namelist.
%
Probably need to also add error checking on if both methods are \quoterm{active}
and whether that is intentional or not (or do we just use the above
precedence?).
%
} % end devNotes
%
\vspace{1.0em} % Add vertical space before next item.

%==========================================================================

%%% The following notes should help when updating this manual in the future.
%%%
%%% For cases that support an internally generated grid (IG),
%%% see subroutine create_internal_grid in internal_grid_mod.f90.
%%% itestcase:
%%%   |Transport_Problems|   = |[Vortex_Problems,Advec_Problems]|
%%%                          = |[Shu_Vortex,Vortex_Transport,Shu_MultiVortex,
%%%                              Density_Transport,Entropy_Transport]|
%%%                          = |[5,6,55,7,8]|
%%%    Taylor_Green_Vortex   =   9
%%%   |Taylor_Couette_Flow|  = |21|
%%%   |Freestream_TCF|       = |22|
%%%   |Pipe_Channel_Flow|    = |25|
%%%   |CH4_Flame_Speed_Test| = |19|
%%%   |Sod_Shock_Tube|       = |23|
%%%   |Density_Pulse_Test|   = |24|
%%%   |Shu_Osher_Test|       = |26|
%%%   |Advection_JiangShu|   = |27|
%%%
%%% For cases that modify the reference conditions (RC),
%%% see subroutine initialize_reference in initialization_mod.f90.
%%% itestcase:
%%%   Advec_Problems = [Density_Transport, Entropy_Transport] = [7,8]
%%%  |Shu_Vortex|         =  5,-5
%%%  |Shu_MultiVortex|    =  55,-55
%%%  -Vortex_Transport    = -6
%%%   Taylor_Green_Vortex =  9
%%%   Taylor_Couette_Flow = 21
%%%   Sod_Shock_Tube      = 23
%%%   Shu_Osher_Test      = 26
%%%   Advection_JiangShu  = 27
%%%
%%% For cases that provide an exact solution (ES),
%%% see sub output_variable in write_cgns_mod.f90.
%%%   Sod_Shock_Tube         = 23, ExactDensity, ExactPressure
%%%   Advection_JiangShu     = 27, ExactDensity, ExactPressure
%%%
%%% FIXME:
%%% For cases that modiy the post-procesing (PP),
%%% see sub init_generic_mod in generic_mod.f90 that writes screen info.
%%% itestcase:
%%%   Inviscid_Gaussian_Bump =  3
%%%   Laminar_BndLyr         =  4
%%%   THX_Phase_3            = 20
%%% see sub get_output_variable in postproc_mod.f90.
%%% itestcase:
%%%   Inviscid_Gaussian_Bump =  3, Entropy calc

\begin{nmldescitem}
Initialization = 1
\end{nmldescitem}
\inp*[nmlInitialization::Initialization]{}[Initialization][]

\begin{nmldescbody}
Specifies the test-case\slash problem to solve, primarily affecting the
initialization and post-processing of the simulation.

\begingroup % limit scope of local commands
\setlength{\tabcolsep}{0.75\tabcolsep} % shrink column gap
%%% If descriptions are too long, need p{} to line wrap.
%\begin{longtable}{ @{} r c c c p{0.75\linewidth} @{} }
\begin{longtable}{ @{} r c c c l @{} }
 \toprule
 Value & IG & RC & PP & Description \\
 \midrule
 \endhead
    1 & . & . & . & Generic flow, initialize to reference conditions.\\
    2 & Y & . & . & Channel flow.\\
    3 & . & . & Y & Inviscid Gaussian bump.\\
    4 & . & . & Y & Laminar subsonic boundary layer.\\
    5 & Y & Y & . & Shu isentropic Euler vortex, propagating diagonally.\\
                    \nopagebreak
   -5 & Y & Y & . & Shu isentropic Euler vortex, stationary.\\
   55 & Y & Y & . & Shu isentropic Euler vortex, propagating diagonally,
                    4 total vortices.\\
                    \nopagebreak
  -55 & Y & Y & . & Shu isentropic Euler vortex, stationary,
                    4 total vortices.\\
    6 & Y & . & . & Isentropic Euler vortex propagating in
                    \inp[nmlReference::AoAAlpha]{AoA\_Alpha}
                    direction.\\
                    \nopagebreak
   -6 & Y & Y & . & Isentropic Euler vortex propagating only in the
                    $y$-direction.\\
    7 & Y & Y & . & Linear advection equation advecting a density profile.\\
    8 & Y & Y & . & Linear advection equation advecting an entropy profile.\\
    9 & Y & Y & . & Taylor-Green vortex problem.\\
   10 & . & . & . & Subsonic jet cases using the SMC000 nozzle.\\
                    \nopagebreak
  -10 & . & . & . & Subsonic jet cases using the ARN2 nozzle.\\
   11 & . & . & . & Supersonic jet cases.\\
   12 & . & . & . & Infinite cylinder.\\
   13 & . & . & . & Freestream preservation.\\
   14 & . & . & . & Reacting backstep.\\
   15 & . & . & . & 2D diffusion flame fueled with hydrogen.\\
   16 & . & . & . & 2D diffusion flame fueled with methane.\\
   17 & . & . & . & Nitrogen\slash Oxygen diffusion.\\
   18 & . & . & . & Compression corner.\\
   19 & Y & . & . & Methane flame speed test.\\
   20 & . & . & Y & THX Phase 3.\\
   21 & Y & Y & . & Taylor-Couette flow.\\
   22 & Y & . & . & Freestream Taylor-Couette flow.\\
   23 & Y & Y & . & 1D Sod shocktube problem.\\
   24 & Y & . & . & 1D Density pulse test.\\
   25 & Y & . & . & Annular pipe channel flow.\\
   26 & Y & Y & . & 1D Shu-Osher problem.\\
   27 & Y & Y & . & 1D Jiang-Shu density pulse test.\\
   \bottomrule
%  \multicolumn{2}{l}{%
%    \footnotesize{%
%      $^\dagger$
%      Supports 
%      \inp[nmlGrid::GridFormat]{Grid\_Format} = -1 (internally generated grid),
%      but external grids can be used as well.
%    }
%  }
\end{longtable}
\endgroup

\noindent
In the list above, the following designations are made:
\begin{enumerate}[align=parleft]
\item[IG]
This input option supports an
\key*[key:InternallyGeneratedGrid]{internally generated grid}[=]
by setting \inp[nmlGrid::GridFormat]{Grid\_Format} = -1 (internal grid),
but external grids can be used as well.
Grid resolution can be controlled as described under the
\inp[nmlGrid::InternalGridDoF]{Internal\_Grid\_DoF} input.
\item[RC]
This input option overrides any user-specified reference conditions
in \inp[nmlReference]{nml\_reference} with preset values inside the code,
which will be written to the screen during code execution.
\item[PP]
This input option only affects certain post-processed quantities,
such as the screen output,
and the flow will be initialized the same as option 1 (generic flow).
\end{enumerate}

\NOTE
This initialization is superseded when either loading a restart file by
specifying \inp[nmlInitialization::LoadRestartFile]{Load\_Restart\_File}=\texttt{T}
or using the input \inp[nmlInitialization::InitialSolutionFile]{Initial\_Solution\_File}
to initializing a simulation from a valid solution file.
%
Instead, the flowfield will be read from the corresponding file.
%
See the description of the
\key[key:InitializationSequence]{initialization sequence}
on \pageref{key:InitializationSequence}.

\end{nmldescbody}

%==========================================================================

\begin{nmldescitem}
Jet_Init_Func = 2
\end{nmldescitem}
\inp*[nmlInitialization::JetInitFunc]{}[Jet\_Init\_Func][]

\begin{nmldescbody}
Sometimes when initializing a jet simulation from freestream conditions,
early transients within the nozzle can be numerically unstable,
and a better initial solution is desirable.

If \inp[nmlInitialization::Initialization]{Initialization}
is $\pm 10$ or $11$, \inptext{Jet\_Init\_Func} can be used to
specify different error function distributions to initialize the
solution for the ARN2, SMC000, and the supersonic jet geometries.
%
The jet geometry (\forexample ARN2, SMC000\cite{Bridges_2011}) used for the
initialization function depends on the value of \inptext{Initialization}.

\begin{tabular}{ @{} r @{ = } p{0.85\linewidth} @{} }
  1 & Constant radial blending. \\
  2 & Blending in both x and r with a constant radial spreading. \\
  3 & Blending in both x and r with a downstream radial spreading.
\end{tabular}
\\
For any other value, the volume will be initialized to the
freestream reference conditions.

\IDEA{More detail could be provided on the profile functions.}
\end{nmldescbody}

%==========================================================================

\begin{nmldescitem}
Load_Restart_File = F
\end{nmldescitem}
\inp*[nmlInitialization::LoadRestartFile]{}[Load\_Restart\_File][]

\begin{nmldescbody}
Specifies whether or not to read a restart file.

\begin{tabular}{ @{} r @{ = } p{0.85\linewidth} @{} }
  T & Read and initialize from the restart file given by 
      \inp[nmlInitialization::RestartFile]{Restart\_File}.\\
  F & Do not read a restart file, but instead initialize according to
      \inp[nmlInitialization::Initialization]{Initialization}.
\end{tabular}

\NOTE
This will be superseded by
\inp[nmlInitialization::InitialSolutionFile]{Initial\_Solution\_File} if the
  specified value is an existing file and
\inp[nmlInitialization::InitialSolutionFormat]{Initial\_Solution\_Format} is set to a valid
value.
%
The only difference is that this option keeps the iteration number and
simulation time from the restart file, whereas \inptext{Initial\_Solution\_File}
resets both to 0.
%
See the description of the
\key[key:InitializationSequence]{initialization sequence}
on \pageref{key:InitializationSequence}.
\end{nmldescbody}

%==========================================================================

\begin{nmldescitem}
Restart_File = "out.rst"
\end{nmldescitem}
\inp*[nmlInitialization::RestartFile]{}[Restart\_File][]

\begin{nmldescbody}
File path and name to the solution restart file that will be \emph{read} if
\inp[nmlInitialization::LoadRestartFile]{Load\_Restart\_File}
is true.
%
%This takes precedence over
This supersedes
\inp[nmlInitialization::Initialization]{Initialization},
but will be superseded by
\inp[nmlInitialization::InitialSolutionFile]{Initial\_Solution\_File}.
%
See the description of the
\key[key:InitializationSequence]{initialization sequence}
on \pageref{key:InitializationSequence}.
%
The solution will continue from the iteration count and simulation time
of the restart file.

\NOTE
This has no effect on the restart files that are \emph{written},
typically to:
\\ % Prevent filename from entering right margin.
  \texttt{Restart\_files/out.<iteration>.rst}.
\end{nmldescbody}

%==========================================================================

%\begin{nmldescitem}
%CFL_Reduction_For_Nan_Restart = 2.0
%\end{nmldescitem}
%\inp*[nmlInitialization::CFLReductionForNanRestart]{}[CFL\_Reduction\_For\_Nan\_Restart][]

%\begin{nmldescbody}
%\NOTE
%This feature is not working.
%
%This feature is intended to provide for an automatic recovery
%for failed solutions by rolling back to a previous state
%and reducing the timestep.
%%
%If an NaN is encountered in a solution, the code will restart from
%the most recent restart file.
%%
%The CFL from when the NaN occurred is divided by
%\inptext{CFL\_Reduction\_For\_Nan\_Restart} and used to continue
%the simulation.
%
%This feature is most useful on computing systems where manually
%modifying and resubmitting a job incurs long delays due to the
%queuing system.
%%
%Unfortunately, in many instances this technique is insufficient
%to overcome the numerical exception, and manual adjustments
%to the limiting \andor filtering are necessary.
%
%%\FIXME{
%%\begin{enumerate}
%%\item
%%REMOVE this input from the code and documentation. Does not work well.
%%\item
%%Current code bypasses automatic restart if
%%  $\inptext{CFL\_Reduction\_For\_Nan\_Restart} > 1$.
%%\item
%%There is no other input that enables\slash disables this automatic recovery.
%%  It is\slash was intended to always be on.
%%\item
%%What happens if the input file uses CFL ramping?  Does ramping continue?
%%\end{enumerate}
%%}
%\end{nmldescbody}

%==========================================================================

\begin{nmldescitem}
Initial_Solution_File = ""
\end{nmldescitem}
\inp*[nmlInitialization::InitialSolutionFile]{}[Initial\_Solution\_File][]

\begin{nmldescbody}
Path to a file that can be \emph{read} for initializing the solution.
%
This can be an externally generated solution or a GFR restart file.
%
This is essentially the same as loading a restart file,
except the \emph{\textbf{iteration count and simulation time are reset to 0}}.
%
See \key[key:InitializationSequence]{initialization sequence}
on \pref{\pageref{key:InitializationSequence}} for important details.
%
When specifying \inptext{Initial\_Solution\_File}, one must also specify a valid value
for \inp[nmlInitialization::InitialSolutionFormat]{Initial\_Solution\_Format} as its default
value will prevent this file from being read.

\NOTE
This has no effect on the solution files that are \emph{written}.

\NOTE
After initialization is complete, set \inptext{Initial\_Solution\_File}=\texttt{""}
so that subsequent simulations will continue from the previous restart file
rather than starting over again.
See \pref{\pageref{key:InitializationSequence}}.


\end{nmldescbody}

%==========================================================================

\begin{nmldescitem}
Initial_Solution_Format = -99  ! -99 = Unknown
\end{nmldescitem}
\inp*[nmlInitialization::InitialSolutionFormat]{}[Initial\_Solution\_Format][]

\begin{nmldescbody}
Format of the file specified by input
\inp[nmlInitialization::InitialSolutionFile]{Initial\_Solution\_File}.

\begin{tabular}{ @{} r @{ = } p{0.85\linewidth} @{} }
  0 & Restart file format.\\
  3 & Plot3D file format (experimental).
\end{tabular}
\\
\NOTE
Any other value will prevent the file specified by
\inp[nmlInitialization::InitialSolutionFile]{Initial\_Solution\_File}
from being read, and another means will be used to initialize the solution.

\end{nmldescbody}

%==========================================================================

